\section{Introduction}
\label{sec:intro}

Large Language Models (LLMs) are increasingly proposed as tools for legal reasoning tasks, from statutory question answering\cite{louis2024interpretable} to case outcome prediction\cite{jiang2023legal}. Yet the same properties that make LLMs attractive, which is its fluent, free-form generation over unstructured text, also make them poorly suited to the standards the legal domain demands\cite{mik2024caveatlectorlargelanguage}. A legal decision is not merely a label, it is expected to be traceable to the rule that produced it, defensible against a specific line of argument, and reproducible under scrutiny. Chain-of-thought prompting improves the appearance of reasoning\cite{yu2023exploring}, but the underlying inference remains implicit, and free-form justifications frequently contain unsupported logical leaps that are difficult to detect from the output text alone\cite{galitsky2026discourse}. For high-stakes legal applications, opacity of this kind is not a minor limitation but a barrier to adoption.

Neuro-symbolic architectures have been proposed as a response to this problem: an LLM is used to extract structured, human-interpretable facts (predicates) from a natural-language fact pattern, and a symbolic reasoner then applies an explicit, auditable rule structure to those facts to reach a decision\cite{yang2025neuro}. Because the rule structure is fixed and inspectable, the resulting decision path can, in principle, be verified independently of the model that produced it. In this paper, we instantiate this division of labour as PYTHEN, a neuro-symbolic pipeline for hearsay classification under the LegalBench benchmark\cite{guha2023legalbench}, and show that it is empirically beneficial: replacing an LLM's final-answer judgment with a symbolic rule tree driven by LLM-extracted predicates improves F1-score and reduces instability across model families, relative to both direct prompting and chain-of-thought prompting. An oracle experiment, in which gold predicates are fed directly into the rule tree, further shows that the symbolic component reaches perfect agreement with ground truth once its inputs are correct - indicating that essentially all of the residual error in the neuro-symbolic pipeline originates upstream, in the LLM's predicate extraction step.

This result, however, exposes a gap that standard evaluation practice does not close. Prior work in relation extraction and structured information extraction commonly evaluates systems using corpus-level aggregate measures, particularly precision, recall, and F1 variants\cite{olek2023evaluationf1scorerecent,chen2023unifiedviewevaluationmetrics}. Such metrics summarize performance over extracted structures or relation instances, rather than exposing quality differences for individual predicate slots. This assumption does not hold once predicates are consumed by a rule tree. A rule tree encodes conditional structure: some predicates sit on a decisive branch where a single incorrect value flips the final verdict, while others are redundant with correlated conditions, are short-circuited by an earlier branch, or belong to a path the specific fact pattern never reaches. An extraction error on the first kind of predicate is decision-critical; the same kind of error on the second kind of predicate is, for that instance, functionally inert. Two models with identical predicate-level F1-scores can therefore differ substantially in how many of their errors actually reach the final decision, and a model that appears weaker by aggregate F1 may in practice make fewer decision-altering mistakes. Aggregate predicate metrics cannot distinguish these cases, because they are computed independently of the rule structure the predicates feed into.

This paper addresses that gap by asking, for a neuro-symbolic legal reasoning pipeline, not only \emph{how often} an LLM extracts a predicate incorrectly, but \emph{whether that error matters}, that is, whether correcting it would change the system's output. We introduce a predicate-level intervention methodology: for each fact pattern, we take the LLM-extracted predicate set, correct one erroneous predicate at a time toward its gold value, and rerun the PYTHEN rule tree to observe whether the decision changes. This yields, for every predicate error observed in practice, a direct measurement of its downstream decision impact, which we compare against the predicate's standalone extraction F1-score. We further refine an extraction-error taxonomy for this pipeline by adding a rule-position dimension, distinguishing decision-critical extraction errors from decision-irrelevant ones, and use it to characterize where LLM extraction failures actually threaten the reliability of the final legal conclusion. The hearsay classification task from LegalBench\cite{guha2023legalbench}, with its validated PYTHEN rule tree, provides a controlled setting in which every predicate's position in the rule structure is known, making such an attribution analysis tractable.

The contribution of this paper is therefore a methodology and an empirical analysis for rule-aware error attribution in neuro-symbolic legal reasoning, rather than a new architecture or a new headline accuracy result. Concretely, we (i) propose a predicate-level intervention procedure for measuring the downstream decision impact of individual predicate extraction errors in a symbolic legal engine; (ii) quantify, for the hearsay classification task, the extent to which predicate-level F1 over- or under-states a predicate's practical importance to the final decision; and (iii) refine the existing extraction error taxonomy to incorporate this decision-impact dimension, offering a more actionable diagnostic for where LLM-symbolic pipelines should be hardened. In doing so, we argue that evaluating neuro-symbolic legal systems requires looking beyond predicate accuracy in isolation, toward how errors propagate through the specific rule structure that governs the decision.

The remainder of this paper is organized as follows. Section~\ref{sec:related} situates this work relative to symbolic and hybrid approaches to legal reasoning and to error analysis in neuro-symbolic systems. Section~\ref{sec:method} formalizes the PYTHEN rule tree, the predicate-level intervention procedure, and the refined error taxonomy. Section~\ref{sec:experiments} reports the resulting decision-impact measurements on the hearsay classification task and compares them against standalone predicate F1. Section~\ref{sec:discussion} discusses the implications for evaluating and hardening neuro-symbolic legal pipelines, and Section~\ref{sec:conclusion} concludes.
